\section{Weekly Summary}

\begin{tcolorbox}[colback=yellow!10,colframe=black!50,title=AI-Generated Content]
\noindent The summary below was written by Claude (Anthropic) from that week's regression outputs and series descriptions. It has not been reviewed or fact-checked by a human, and should be read as an automated, informal recap -- not analysis.
\end{tcolorbox}

Welcome back to the corner of the internet where we throw two unrelated economic series into a blender every day and act surprised when soup comes out. This week's contestants included a heavy rotation of California house prices, a couple of discontinued banking indicators that FRED clearly keeps around out of sentimentality, and — for spice — Japanese retail inventory sentiment and the carbon dioxide emissions of commercial-sector distillate fuel. Truly the crossover episodes nobody demanded.

We kicked things off Monday with Oxnard-Thousand Oaks-Ventura house prices versus the gross fixed investment of financial businesses in equipment, software, and structures, a series so discontinued it stopped reporting in 2017 and still got dragged into this. The raw regression posted an R-squared of 0.87 and a p-value with more zeros after the decimal than I have patience to count, which is the statistical equivalent of a fireworks show. Except, of course, it's two things that both went up for fifty years. Once we detrended and stripped out that shared "everything drifts upward over time" magic, the R-squared face-planted to 0.23. Still technically significant, sure, but it went from "soulmates" to "we follow each other on LinkedIn." Classic trend mirage, exactly the kind of thing this project exists to point and laugh at.

Tuesday brought San Luis Obispo house prices against the percent of loans that are prime-based by average loan size, and honestly this one had the decency to flop right out of the gate. The raw regression couldn't even clear the significance bar (p around 0.076, R-squared a majestic 0.04), and detrending only made it sadder. No trend mirage, no genuine link, just two series politely ignoring each other. Refreshing, in a way. Then Wednesday served up commercial distillate fuel CO2 emissions versus Japanese retail stock-volume sentiment, which is where it gets mildly interesting: the raw regression showed a positive relationship (R-squared 0.48), but after detrending the sign flipped to negative and it stayed significant at the 0.017 level. So the "real" relationship, insofar as such a phrase is legally permissible here, is the opposite of what the raw version claimed. Dramatic! Meaningful? Absolutely not.

As always, the mandatory reminder: these are randomly paired series with zero causal story and no rigorous statistical claim being made, so please do not go rebalancing your portfolio based on Japanese inventory vibes and diesel fumes. Spurious correlation isn't the occasional bug in this project — it's the entire product. This week the detrending step did its usual job of quietly assassinating one over-hyped raw "finding" and gently roughing up another, which is roughly as much scientific integrity as you can expect from a machine that thinks Ventura County real estate and discontinued financial-investment flows are destiny. See you next week for more of the same nonsense.
